\documentclass[12pt]{article}
\usepackage{amsmath, amssymb, amsthm, enumerate, graphicx}
\usepackage[usenames,dvipsnames]{color}
\usepackage{bm}
\usepackage[colorlinks=true,urlcolor=blue]{hyperref}
\usepackage{geometry}
\geometry{margin=1in}
\usepackage{float}
\usepackage{graphics}
\setlength{\marginparwidth}{2.15cm}
\usepackage{booktabs}
\usepackage{enumitem}
\usepackage{epsfig}
\usepackage{setspace}
\usepackage{parskip}
\usepackage[normalem]{ulem}
\usepackage{tikz}
\usetikzlibrary{positioning, arrows, automata}
\usepackage{pgfplots}
\usepackage[font=scriptsize]{subcaption}
\usepackage{float}
\usepackage[]{algorithm2e}
\usepackage{environ}
\usepackage{bbm}
\usepackage{titling}
\usepackage{url}
\usepackage{xcolor}
\usepackage{lipsum}
\usepackage{lastpage}
\usepackage[colorlinks=true,urlcolor=blue]{hyperref}
\usepackage{multicol}
\usepackage{tabularx}
\usepackage{comment}
\usepackage[utf8]{inputenc}
\usepackage{amssymb}
\usepackage{setspace}
\usepackage{marvosym}
\usepackage{wrapfig}
\usepackage{datetime}
\usepackage[many]{tcolorbox}
\usepackage{array}
\usepackage{multirow}
\usepackage{wasysym}
\usepackage{cancel}
\usepackage{fancyhdr}
\usepackage{listings}
\usepackage{color}
\usepackage{soul}







%%% old packages
\usepackage{fullpage}
\usepackage{sectsty}
\usepackage{amsfonts}


\newcommand{\R}{\mathbb{R}}
\newcommand{\blackcircle}{\tikz\draw[black,fill=black] (0,0) circle (1ex);}
\renewcommand{\circle}{\tikz\draw[black] (0,0) circle (1ex);}
\DeclareMathOperator{\sign}{sign}


%%%%%%%%%%%%%%%%%%%%%%%%%%%%%%%%%%%%%%%%%%
% Custom commands                        %
%%%%%%%%%%%%%%%%%%%%%%%%%%%%%%%%%%%%%%%%%%

\newcommand{\vc}[1]{\boldsymbol{#1}}
\newcommand{\adj}[1]{\frac{d J}{d #1}}
\newcommand{\chain}[2]{\adj{#2} = \adj{#1}\frac{d #1}{d #2}}

% mathcal
\newcommand{\Ac}{\mathcal{A}}
\newcommand{\Bc}{\mathcal{B}}
\newcommand{\Cc}{\mathcal{C}}
\newcommand{\Dc}{\mathcal{D}}
\newcommand{\Ec}{\mathcal{E}}
\newcommand{\Fc}{\mathcal{F}}
\newcommand{\Gc}{\mathcal{G}}
\newcommand{\Hc}{\mathcal{H}}
\newcommand{\Ic}{\mathcal{I}}
\newcommand{\Jc}{\mathcal{J}}
\newcommand{\Kc}{\mathcal{K}}
\newcommand{\Lc}{\mathcal{L}}
\newcommand{\Mc}{\mathcal{M}}
\newcommand{\Nc}{\mathcal{N}}
\newcommand{\Oc}{\mathcal{O}}
\newcommand{\Pc}{\mathcal{P}}
\newcommand{\Qc}{\mathcal{Q}}
\newcommand{\Rc}{\mathcal{R}}
\newcommand{\Sc}{\mathcal{S}}
\newcommand{\Tc}{\mathcal{T}}
\newcommand{\Uc}{\mathcal{U}}
\newcommand{\Vc}{\mathcal{V}}
\newcommand{\Wc}{\mathcal{W}}
\newcommand{\Xc}{\mathcal{X}}
\newcommand{\Yc}{\mathcal{Y}}
\newcommand{\Zc}{\mathcal{Z}}

% mathbb
\newcommand{\Ab}{\mathbb{A}}
\newcommand{\Bb}{\mathbb{B}}
\newcommand{\Cb}{\mathbb{C}}
\newcommand{\Db}{\mathbb{D}}
\newcommand{\Eb}{\mathbb{E}}
\newcommand{\Fb}{\mathbb{F}}
\newcommand{\Gb}{\mathbb{G}}
\newcommand{\Hb}{\mathbb{H}}
\newcommand{\Ib}{\mathbb{I}}
\newcommand{\Jb}{\mathbb{J}}
\newcommand{\Kb}{\mathbb{K}}
\newcommand{\Lb}{\mathbb{L}}
\newcommand{\Mb}{\mathbb{M}}
\newcommand{\Nb}{\mathbb{N}}
\newcommand{\Ob}{\mathbb{O}}
\newcommand{\Pb}{\mathbb{P}}
\newcommand{\Qb}{\mathbb{Q}}
\newcommand{\Rb}{\mathbb{R}}
\newcommand{\Sb}{\mathbb{S}}
\newcommand{\Tb}{\mathbb{T}}
\newcommand{\Ub}{\mathbb{U}}
\newcommand{\Vb}{\mathbb{V}}
\newcommand{\Wb}{\mathbb{W}}
\newcommand{\Xb}{\mathbb{X}}
\newcommand{\Yb}{\mathbb{Y}}
\newcommand{\Zb}{\mathbb{Z}}

% mathbf lowercase
\newcommand{\av}{\mathbf{a}}
\newcommand{\bv}{\mathbf{b}}
\newcommand{\cv}{\mathbf{c}}
\newcommand{\dv}{\mathbf{d}}
\newcommand{\ev}{\mathbf{e}}
\newcommand{\fv}{\mathbf{f}}
\newcommand{\gv}{\mathbf{g}}
\newcommand{\hv}{\mathbf{h}}
\newcommand{\iv}{\mathbf{i}}
\newcommand{\jv}{\mathbf{j}}
\newcommand{\kv}{\mathbf{k}}
\newcommand{\lv}{\mathbf{l}}
\newcommand{\mv}{\mathbf{m}}
\newcommand{\nv}{\mathbf{n}}
\newcommand{\ov}{\mathbf{o}}
\newcommand{\pv}{\mathbf{p}}
\newcommand{\qv}{\mathbf{q}}
\newcommand{\rv}{\mathbf{r}}
\newcommand{\sv}{\mathbf{s}}
\newcommand{\tv}{\mathbf{t}}
\newcommand{\uv}{\mathbf{u}}
\newcommand{\vv}{\mathbf{v}}
\newcommand{\wv}{\mathbf{w}}
\newcommand{\xv}{\mathbf{x}}
\newcommand{\yv}{\mathbf{y}}
\newcommand{\zv}{\mathbf{z}}

% mathbf uppercase
\newcommand{\Av}{\mathbf{A}}
\newcommand{\Bv}{\mathbf{B}}
\newcommand{\Cv}{\mathbf{C}}
\newcommand{\Dv}{\mathbf{D}}
\newcommand{\Ev}{\mathbf{E}}
\newcommand{\Fv}{\mathbf{F}}
\newcommand{\Gv}{\mathbf{G}}
\newcommand{\Hv}{\mathbf{H}}
\newcommand{\Iv}{\mathbf{I}}
\newcommand{\Jv}{\mathbf{J}}
\newcommand{\Kv}{\mathbf{K}}
\newcommand{\Lv}{\mathbf{L}}
\newcommand{\Mv}{\mathbf{M}}
\newcommand{\Nv}{\mathbf{N}}
\newcommand{\Ov}{\mathbf{O}}
\newcommand{\Pv}{\mathbf{P}}
\newcommand{\Qv}{\mathbf{Q}}
\newcommand{\Rv}{\mathbf{R}}
\newcommand{\Sv}{\mathbf{S}}
\newcommand{\Tv}{\mathbf{T}}
\newcommand{\Uv}{\mathbf{U}}
\newcommand{\Vv}{\mathbf{V}}
\newcommand{\Wv}{\mathbf{W}}
\newcommand{\Xv}{\mathbf{X}}
\newcommand{\Yv}{\mathbf{Y}}
\newcommand{\Zv}{\mathbf{Z}}

% bold greek lowercase
\newcommand{\alphav     }{\boldsymbol \alpha     }
\newcommand{\betav      }{\boldsymbol \beta      }
\newcommand{\gammav     }{\boldsymbol \gamma     }
\newcommand{\deltav     }{\boldsymbol \delta     }
\newcommand{\epsilonv   }{\boldsymbol \epsilon   }
\newcommand{\varepsilonv}{\boldsymbol \varepsilon}
\newcommand{\zetav      }{\boldsymbol \zeta      }
\newcommand{\etav       }{\boldsymbol \eta       }
\newcommand{\thetav     }{\boldsymbol \theta     }
\newcommand{\varthetav  }{\boldsymbol \vartheta  }
\newcommand{\iotav      }{\boldsymbol \iota      }
\newcommand{\kappav     }{\boldsymbol \kappa     }
\newcommand{\varkappav  }{\boldsymbol \varkappa  }
\newcommand{\lambdav    }{\boldsymbol \lambda    }
\newcommand{\muv        }{\boldsymbol \mu        }
\newcommand{\nuv        }{\boldsymbol \nu        }
\newcommand{\xiv        }{\boldsymbol \xi        }
\newcommand{\omicronv   }{\boldsymbol \omicron   }
\newcommand{\piv        }{\boldsymbol \pi        }
\newcommand{\varpiv     }{\boldsymbol \varpi     }
\newcommand{\rhov       }{\boldsymbol \rho       }
\newcommand{\varrhov    }{\boldsymbol \varrho    }
\newcommand{\sigmav     }{\boldsymbol \sigma     }
\newcommand{\varsigmav  }{\boldsymbol \varsigma  }
\newcommand{\tauv       }{\boldsymbol \tau       }
\newcommand{\upsilonv   }{\boldsymbol \upsilon   }
\newcommand{\phiv       }{\boldsymbol \phi       }
\newcommand{\varphiv    }{\boldsymbol \varphi    }
\newcommand{\chiv       }{\boldsymbol \chi       }
\newcommand{\psiv       }{\boldsymbol \psi       }
\newcommand{\omegav     }{\boldsymbol \omega     }

% bold greek uppercase
\newcommand{\Gammav     }{\boldsymbol \Gamma     }
\newcommand{\Deltav     }{\boldsymbol \Delta     }
\newcommand{\Thetav     }{\boldsymbol \Theta     }
\newcommand{\Lambdav    }{\boldsymbol \Lambda    }
\newcommand{\Xiv        }{\boldsymbol \Xi        }
\newcommand{\Piv        }{\boldsymbol \Pi        }
\newcommand{\Sigmav     }{\boldsymbol \Sigma     }
\newcommand{\Upsilonv   }{\boldsymbol \Upsilon   }
\newcommand{\Phiv       }{\boldsymbol \Phi       }
\newcommand{\Psiv       }{\boldsymbol \Psi       }
\newcommand{\Omegav     }{\boldsymbol \Omega     }


% SOLUTION environment
\NewEnviron{soln}{
\leavevmode\color{red}\ignorespaces \textbf{Solution} \BODY }{}

% QUESTION AUTHORS environment
\NewEnviron{qauthor}{
\leavevmode\color{blue}\ignorespaces \textbf{Author} \BODY}{}

%%%%%%%%%%%%%%%%%%%%%%%%%%%%%%%%%%%%%%%%%%%%%%%%%%%%%%%%%%%%%%%%%%%%%%%%%%%%%%
%%%%%%%%%%%%%%%%%%%%%%%%%%%%%%%%%%%%%%%%%%%%%%%%%%%%%%%%%%%%%%%%%%%%%%%%%%%%%%
%%%%%%%%%%%%%%%%%%%%%%%%%%%%%%%%%%%%%%%%%%%%%%%%%%%%%%%%%%%%%%%%%%%%%%%%%%%%%%
% TO ONLY SHOW HOMEWORK QUESTIONS, include following (else comment out):
 %\RenewEnviron{soln}{}
 % \RenewEnviron{qauthor}{}
%%%%%%%%%%%%%%%%%%%%%%%%%%%%%%%%%%%%%%%%%%%%%%%%%%%%%%%%%%%%%%%%%%%%%%%%%%%%%%
%%%%%%%%%%%%%%%%%%%%%%%%%%%%%%%%%%%%%%%%%%%%%%%%%%%%%%%%%%%%%%%%%%%%%%%%%%%%%%
%%%%%%%%%%%%%%%%%%%%%%%%%%%%%%%%%%%%%%%%%%%%%%%%%%%%%%%%%%%%%%%%%%%%%%%%%%%%%%

%\newcommand{\norm}[1]{\lVert #1 \rVert}
%\newcommand{\st}{\mathrm{s.t.}}

\makeatletter
\newcommand{\removelatexerror}{\let\@latex@error\@gobble}
\makeatother

\newcommand{\argmax}{\mathop{\mathrm{argmax}}}
\newcommand{\argmin}{\mathop{\mathrm{argmin}}}


%%%%%%%%%%%%%%%%%%%%%%%%%%%%%%%%%%%%%%%%%%%
% Custom box for highlights               %
%%%%%%%%%%%%%%%%%%%%%%%%%%%%%%%%%%%%%%%%%%%

% Define box and box title style
\tikzstyle{mybox} = [fill=blue!10, very thick,
    rectangle, rounded corners, inner sep=1em, inner ysep=1em]

\NewEnviron{notebox}{
\begin{tikzpicture}
\node [mybox] (box){
    \begin{minipage}{\textwidth}
        \BODY
    \end{minipage}
};
\end{tikzpicture}
}


\begin{document}

\section*{}
\begin{center}
  \centerline{\textsc{\LARGE  Homework 2}}
  \vspace{0.5em}
  \centerline{\textsc{\large Linear Regression, Logistic Regression, MLE/MAP, Na\"ive Bayes}
  }
  \vspace{1em}
  \textsc{\large UTK COSC 522: Machine Learning (Fall 2025)} \\
  \vspace{0.5em}

  \centerline{OUT: Sunday, Sep 21st, 2025}
  %\today{} at \currenttime{}}}
  \vspace{0.5em}
  \centerline{\textbf{DUE: Monday, Oct 13th, 2025, 11:59pm}, \textbf{Bonus Point Deadline: Oct 8th, 2025, 11:59pm}}
\end{center}


\section*{Homework Instructions}
\begin{itemize}
\item \textbf{Collaboration policy: Taking Responsibility for Your Career} 

The skills you learn in this course are tools for your future, and this policy is designed to help you sharpen them. You are responsible for your own learning and career, and this framework ensures you get the most from every assignment.

We encourage collaboration, but it must be done \emph{right}. \emph{First}, make a genuine effort to solve problems on your own. This independent effort is where the most critical learning occurs. \emph{Afterward}, you may discuss strategies with peers or consult resources to clarify concepts—the goal is to deepen your own understanding, not to simply get an answer.

\emph{Finally}, like any professional, you must stand behind your own work. Your submitted solution must be written entirely by you, from scratch, without collaborators present. This proves you have truly mastered the material. To maintain academic and professional integrity, you must also cite any person or resource that contributed to your understanding along the way. 

\item\textbf{Late Submission Policy:} See the late homework policy \href{https://docs.google.com/document/d/1lYopNzjn_0AOipRiiwO4h2Ow64seQNeD/edit?usp=sharing&ouid=114000599194691965836&rtpof=true&sd=true}{here}.

\item\textbf{Submitting your work:} 

\begin{itemize}

\item \textbf{Canvas:} 
For this homework, you will submit a \textbf{single .zip file} to Canvas. Please name your file in the format \textbf{YourName\_NetID\_HW2.zip}. This zip file must contain exactly two files: your written solutions in a single PDF and your Python code.

Your \textbf{single PDF file} must contain your answers to all questions, including written problems (proofs, short answers, plots) and the empirical questions from the programming section. You must use the provided homework template. Using the LaTeX version to typeset is strongly recommended, but you may also submit a scan of the template with legible handwriting; \textbf{illegible answers will not be graded}. Please complete all answers within the provided boxes. 

The second file in your zip archive must be your \emph{single} completed programming file, named exactly as \textbf{time\ series.py}. Please ensure your code is runnable, as we may execute it to verify your results.

Regrade requests can be made after homework grades are released. Please be aware that a regrade request allows the TA to review your entire assignment, which may result in points being deducted if new errors are found.

\end{itemize}

\end{itemize}

%For multiple-choice and select-all-that-apply questions, please shade the corresponding box or circle to indicate your answer(s). If you are using the LaTeX template, use the $\blacksquare$ command for shaded boxes and $\blackcircle$ for shaded circles. \textbf{Note}: Do not alter the template in any other way. If a question includes a box for showing your work, you must \textbf{show your work} to receive full credit.

\clearpage
\section{Regularized Linear Regression [10 Points]}

\begin{enumerate}

\item {\bf [10 Points]} Consider the following linear regression model: for each data point in $\mathcal{D} = \{(\bm{x}^{(i)}, y^{(i)})\}_{i=1}^n$,
\begin{align}
    y^{(i)} &= \bm{w}^T \bm{x}^{(i)} + \epsilon\; \text{where}\; y^{(i)}, \epsilon \in \mathbb{R}\; \text{and}\; \bm{w}, \bm{x}^{(i)} \in \mathbb{R}^{d+1} \nonumber
\end{align}
In matrix notation, we can express this linear relationship for all data points as:
\begin{align}
    \bm{y} &= \bm{X}\bm{w} + \bm{\epsilon}\; \text{where}\; \bm{y}, \bm{\epsilon} \in \mathbb{R}^n, \bm{X} \in \mathbb{R}^{n \times (d+1)},\; \text{and}\; \bm{w} \in \mathbb{R}^{d+1} \nonumber
\end{align}
Assuming the residuals are normal and i.i.d.\ ($\bm{\epsilon} \sim \mathcal{N}(\bm{0}, \sigma^2 \bm{I})$), we can write:
\begin{align}
    \bm{y}|\bm{X},\bm{w} &\sim \mathcal{N}(\bm{X}\bm{w}, \sigma^2 \bm{I}) \nonumber
\end{align}

Now assume that we have a Gaussian prior on $\bm{w}$:
\begin{align}
    \bm{w} &\sim \mathcal{N}\left(0,\frac{2\sigma^2}{\lambda}\bm{I}\right) \nonumber
\end{align}
for some fixed $\lambda > 0$. Recall that in ridge regression, the optimal parameter vector is given by:
\begin{align}
\wv^\star = \argmin_{\wv} \|\yv - X\wv\|_2^2 + \lambda \|\wv\|_2^2. \nonumber
\end{align}

Show that the solution to the MAP estimate $\bm{w}^*_{\text{MAP}}$ in this setting is the same as the one obtained from ridge regression. 

(Hint: Start by writing down the expression for the negative log posterior and show that minimizing it gives the same solution as minimizing the OLS with Ridge regression.)

\begin{tcolorbox}[fit,height=9cm, width=0.95\textwidth, blank, borderline={1pt}{-2pt}]
%solution
\end{tcolorbox}


\end{enumerate}


\clearpage
\section{Convexity of Logistic Regression [28 points]}
Consider a binary classification problem where the goal is to predict a label $y \in \{0,1\}$, given an input $\xv \in \mathbb{R}^d$. A method that you can use for this task is \textit{logistic regression}. In logistic regression, we model the log-odds as an affine function of the data and find weights to maximize the likelihood of our data under the resulting model. 

Recall that an \emph{affine function} $f$ takes the form $f(x)=w^\top x + c$ with $c\in\R$ and $w,x\in\R^n$. In other words, it is a linear function composed with a translation.

\subsection{Convex Optimization}
Recall that the log-likelihood for a logistic regression model can be written as 
\begin{equation*}
    \mathcal{L}(w) = \log P(y|\mathbf{X}, w) = \sum_{i=1}^n[y_i w^\top x_i - \log(1+ \exp(w^\top x_i))].
\end{equation*}
Our goal is to find the weight vector $w$ that maximizes this likelihood. Unfortunately, for this model, we cannot derive a closed-form solution with MLE. An alternative way to solve for $w$ is to use gradient \emph{ascent}, and update $w$ step by step towards the optimal $w$. But we know gradient ascent will converge to the optimal solution $w$ that maximizes the conditional log likelihood $\mathcal{L}$ when $\mathcal{L}$ is concave. In this question, you will prove that $\mathcal{L}$ is indeed a concave function (and hence, the negative conditional log likelihood is a convex function). 

\begin{enumerate}
    \item \textbf{[5 points]} A real-valued function $f:S \rightarrow \mathcal{R}$ defined on a convex set S, is said to be \textit{convex} if $$f(tx_{1}+(1-t)x_{2})\leq tf(x_{1})+(1-t)f(x_{2}), \forall x_{1},x_{2}\in S,\forall t\in [0,1].$$
    Show that a linear combination of $n$ convex functions, $f_1, f_2,...,f_n$, $\sum_{i=1}^n a_if_i(x)$ is also a convex function  $\forall a_i \in R^+$. 


\begin{tcolorbox}[fit,height=7cm,width=.95\textwidth,blank,borderline={1pt}{-2pt},nobeforeafter]
%solution   
\end{tcolorbox}

\clearpage
\item \textbf{[3 point]} Show that a linear combination of $n$ concave functions, $f_1, f_2,...,f_n$, $\sum_{i=1}^n a_if_i(x)$ is also a concave function  $\forall a_i \in R^+$. Recall that if a function $f(x)$ is convex, then $-f(x)$ is concave. (You can use the result from part (1)) 
    

\begin{tcolorbox}[fit,height=7cm,width=.95\textwidth,blank,borderline={1pt}{-2pt},nobeforeafter]
%solution
 \end{tcolorbox}
  
\item \textbf{[5 points]} Another property of twice differentiable convex functions is that the second derivative is non-negative. Using this property, show that $f (x) = \log (1 + \exp x)$ is a convex function. Note that this property is both sufficient and necessary. i.e. (if $f''(x)$ exists, then $f''(x)\ge 0 \iff  f$ is convex ) 
    
\begin{tcolorbox}[fit,height=7cm,width=.95\textwidth,blank,borderline={1pt}{-2pt},nobeforeafter]
%solution
\end{tcolorbox}

\clearpage
\item \textbf{[5 points]} Let $f_i:\mathcal{S} \rightarrow \mathcal{R}$ for $i = 1, \ldots, n$ be a set of convex functions.
Is $f(x) = max_i f_i(x)$ also convex?
If yes, prove it. 
If not, provide a counterexample.  


\begin{tcolorbox}[fit,height=6cm,width=.95\textwidth,blank,borderline={1pt}{-2pt},nobeforeafter]
%solution
\end{tcolorbox}

\item \textbf{[10 points]} Show that the log likelihood of \textit{Logistic Regression} is a concave function. You may use the fact that if $f$ and $g$ are both convex, twice differentiable and $g$ is non-decreasing, then $g \circ f$ is convex. 
    
    
\begin{tcolorbox}[fit,height=12cm,width=.95\textwidth,blank,borderline={1pt}{-2pt},nobeforeafter]
%solution
\end{tcolorbox}
  
\end{enumerate}


\clearpage
\section{Na\"{i}ve  Bayes [32 Points]}

%In this question, we use upper-case letters such as $X, y$ to denote random variables, and lower-case letters to denote \emph{values} of random variables.

Let $X=(x_1,x_2,..,.x_d)$ denote a set of features and $y\in \{0, 1\}$ denote a binary label. Recall that na\"{i}ve Bayes models the conditional label distribution $P(y \mid X)$ via the conditional distribution of features given the label $P(X \mid y)$:
\begin{equation}\label{eqn:pxy}
P(y \mid X) \propto P(X \mid y)P(y) \nonumber
\end{equation}
%\begin{equation}
%P(X|y)P(y).
%\end{equation}

\begin{enumerate}
    \item \textbf{Multinomial Na\"{i}ve Bayes} Suppose that each feature $x_i$ takes values in the set $\{1,2,...,K\}$. Further, suppose that the label distribution is Bernoulli, and the feature distribution conditioned on the label is multinomial. 
    \begin{enumerate}
        \item \textbf{[3 points]} What is the total number of parameters of the model under the na\"{i}ve Bayes assumption? For full credit you must show your work.
    
        \begin{tcolorbox}[fit,height=3cm, width=0.85\textwidth, blank, borderline={1pt}{-2pt}]
        %solution 
        \end{tcolorbox}
        
    \item\textbf{[3 points]} What is the total number of parameters of the model without the na\"{i}ve Bayes assumption? For full credit you must show your work.
    
        \begin{tcolorbox}[fit,height=3cm, width=0.85\textwidth, blank, borderline={1pt}{-2pt}]
        %solution 
        \end{tcolorbox}
        
    \item \textbf{[6 points]} Suppose we change the set of values that $y$ takes, so that $y\in \{0, 1, ..., M-1\}$. How would your answers change in both cases (with and without the na\"{i}ve Bayes assumption)? For full credit you must show your work.
    
        \begin{tcolorbox}[fit,height=5cm, width=0.85\textwidth, blank, borderline={1pt}{-2pt}]
        %solution 
        \end{tcolorbox}
        
    \end{enumerate}
    
    \newpage
    
    \item \textbf{[10 points] Gaussian Na\"{i}ve  Bayes} Now suppose each feature is real-valued, with $x_i \in \mathbb{R}$, and $P(x_i \mid y=c)\sim \mathcal{N}(\mu_{i,c}, 1)$ for $i=1,2,...,d$ and $c=0,1$. Again, suppose that the label distribution is Bernoulli with $P(y = 1) = \pi$. Under the na\"{i}ve Bayes assumption, show that the decision boundary $\{(x_1, x_2, \dots, x_d): P(y=0 \mid x_1, x_2, \dots, x_d) = P(y=1 \mid x_1, x_2, \dots, x_d)\}$ is linear in $x_1, x_2, \dots, x_d$.
    
        \begin{tcolorbox}[fit,height=14cm, width=0.95\textwidth, blank, borderline={1pt}{-2pt}]
        %solution 
        \end{tcolorbox}
        
    \clearpage
    
    \item \textbf{MLE estimators can be biased} Given $N$ independent observations drawn from a univariate Gaussian distribution $x^{(1)}, ..., x^{(N)} \sim \text{N}(\mu, \sigma^2)$, recall that the MLE of the mean and variance parameters are
    \begin{align*}
        \widehat{\mu} = \frac{1}{n}\sum_{i=1}^{n} x^{(i)} \textrm{ and } \widehat{\sigma}^2 = \frac{1}{n} \sum_{i=1}^{n}(x^{(i)} - \widehat{\mu})^2.
    \end{align*}

    \begin{enumerate}
        \item \textbf{[6 points]} Prove that $\widehat{\mu}$ is an \emph{unbiased} estimator of $\mu$ (by showing $\mathbb{E}[\widehat{\mu}] = \mu$) and that $\widehat{\sigma}^2$ is a \emph{biased} estimator of $\sigma^2$ (by showing $\mathbb{E}[\widehat{\sigma}^2] \ne \sigma^2$).
        \begin{tcolorbox}[fit,height=10cm, width=0.85\textwidth, blank, borderline={1pt}{-2pt}]
        % solution
        \end{tcolorbox}

        \item \textbf{[4 points]} Based on your response to the previous question, propose an unbiased estimator of $\sigma^2$.
        \begin{tcolorbox}[fit,height=4cm, width=0.85\textwidth, blank, borderline={1pt}{-2pt}]
        % solution
        \end{tcolorbox}
    \end{enumerate}
\end{enumerate}


\clearpage
\section{Linear Regression for Time Series [30 Points]}

In this problem you will explore fitting a linear regression to predict time series data using OLS and stochastic gradient descent.

\subsection{Data Processing}

In this problem, you will be using a temperature dataset which can be found in the file \texttt{temperature.csv}. This file contains timestamps (column labeled \texttt{`Date Time'}) at 30 minute intervals and corresponding temperature measurements (column labeled \texttt{`T (degC)'}) over eight years of data collection. 

Your first task is to split this data into a train set and a test set. We want to use the first 6 years of data as our training set and the last 2 years as our test set. Since the dataset includes one measurement every 30 minutes (i.e. 2 measurements every hour), the training set should contain the first $2 * 24 * 365 * 6 = 105142$ samples and the test set should contain the last $2 * 24 * 365 * 2 = 35040$ samples (i.e. training and test comprises the entire dataset).

Our first task will be to train a model that predicts the temperature at a given time based on the measurements at the previous \textbf{D=10} timesteps, so we use the value $10$ for $D$ below to set us up for this.

Concretely, we can write our training portion of the dataset as $X_{\mathrm{train}} = \{x_1, x_2, ... ,x_T \}$ where each $x_i$ is one temperature measurement and $T = 105142$ is the total number of training samples. 
In this setting, we will use the temperatures $x_1, x_2, x_3, ..., x_D$ to predict the value at $x_{D+1}$; temperatures $x_2, x_3, ..., x_{D+1}$ to predict the temperature at $x_{D+2}$; and so on so that in general we use $x_i, x_{i+1}, ..., x_{i+(D-1)}$ to predict the value of $x_{i+D}$.

You need to reformat the training portion of the dataset to follow this framework. You should create an \texttt{X\_train} matrix that has $D$ columns (i.e. the $D$ consecutive timestamps) and $T-D$ rows. Note that data will be repeated in this matrix, since each row is shifted by only one timestep from the previous row and will therefore contain $D-1$ of the same temperature values. You should also create a \texttt{y\_train} vector that contains the target values we want to predict, i.e. $[x_{D+1}, x_{D+2}, ..., x_T]$.

Similarly, we can define our test portion of our dataset as $X_{\mathrm{test}} = \{x_{T+1}, x_{T+2}, ... ,x_{T+C} \}$ where $C$ is the total number of the test samples. We will create an \texttt{X\_test} matrix and a \texttt{y\_test} vector following the same procedure as for the training dataset. 


We have now created normalized datasets with 10 features and can use these 10 features to predict the target corresponding to each row. 

\clearpage

\subsection{Predicting Temperatures using Linear Regression}

For this question, please submit all code you wrote in a single, self-contained file titled \texttt{time\_series.py} on Canvas.

\begin{enumerate}
    \item \textbf{[4 points]} Fit an OLS linear regression model using \texttt{X\_train} and \texttt{y\_train} to find the OLS solution. Report the following values:
    \begin{enumerate}
        \item the weights you learned (including the bias or intercept weight)
        \item the time taken to fit the model
        \item the MSE on \texttt{X\_test}
    \end{enumerate}
    As a reminder, you are not permitted to use libraries other than \texttt{numpy} in your implementations. 
    
    \begin{tcolorbox}[fit,height=5cm, width=0.95\textwidth, blank, borderline={1pt}{-2pt}]
    %solution
    \end{tcolorbox}
    
    \item \textbf{[8 points]} Repeat the previous question for $D=\{50, 100, 500\}$. For these models, you should just report the weights on the first 10 features and the bias weight, along with the time required to fit each model and the MSE on \texttt{X\_test}.
    \begin{tcolorbox}[fit,height=7cm, width=0.95\textwidth, blank, borderline={1pt}{-2pt}]
    %solution
    \end{tcolorbox}

    \clearpage
    
    \item \textbf{[3 points]} Using the times taken to fit the models for $D=10, 50, 100,$ and $500$, estimate the time it would take to fit a model using features from an entire year's worth of data (i.e. $D = 2 * 24 * 365 = 17520$). Report how you estimated the time (i.e. what function did you fit?) and your time estimate in minutes.
    \begin{tcolorbox}[fit,height=3cm, width=0.95\textwidth, blank, borderline={1pt}{-2pt}]
    %solution
    \end{tcolorbox}

    \clearpage
    
    \item \textbf{[10 points]} Based on our result from the previous part, we may conclude that using OLS will result in a very high computational cost. As an alternative, you will now learn the weights $\mathbf{w}$ and bias $b$ using \textbf{stochastic gradient descent}. We will use train and test datasets created for $D=17520$, i.e. using the data from an entire year.
    
    Your implementation of SGD should use a learning rate of $\eta=1e-10$ and run for 20 epochs. Initialize your weights to be uniformly distributed, with each value set to $\frac{1}{D}$ and your bias to be $1$. Additionally, while normally you would shuffle or permute the training data points in each epoch of SGD, for this assignment, you should loop through the training dataset in chronological order (i.e., the original ordering of the dataset) without randomly shuffling the data points. 
    
    Again, please only use \texttt{numpy} to implement SGD from scratch. After using SGD to train a LR model that uses a year's worth of data for 20 epochs, please report the following values:
    \begin{enumerate}
        \item Train MSE
        \item Test MSE
        \item Total training time
        \item First 10 items in your final weight vector. 
    \end{enumerate}
    

    \begin{tcolorbox}[fit,height=8cm, width=0.95\textwidth, blank, borderline={1pt}{-2pt}]
    %solution
    \end{tcolorbox}

    \clearpage
    
    \item \textbf{[5 points]} Now we can examine the weights learned by SGD. What takeaways can you observe? Are any features particularly informative in predicting the targets? Please give a 2-3 sentence response describing your observations. \textbf{Hint:} consider which features have the largest weights; what observations do those features correspond to? 
    \begin{tcolorbox}[fit,height=5cm, width=0.95\textwidth, blank, borderline={1pt}{-2pt}]
    %%% Write your solution here %%%
    \end{tcolorbox}
    
\end{enumerate}


\clearpage
\section{Collaboration Questions}

\begin{enumerate}
    \item 
    \begin{enumerate}
        \item Did you receive any help whatsoever from anyone in solving this assignment? 
        \begin{soln} Yes / No. \end{soln} 
        \item If you answered `yes', give full details (e.g. “Jane Doe explained to me what is asked in Question 3.4”)\\
        
        \begin{tcolorbox}[fit,height=4cm, width=15cm, blank, borderline={1pt}{-2pt},nobeforeafter]
        
        \begin{soln}
        \end{soln}
        
        \end{tcolorbox}
        
    \end{enumerate}
    
    
    \item 
    \begin{enumerate}
        \item Did you give any help whatsoever to anyone in solving this assignment? 
        \begin{soln} Yes / No. \end{soln} 
        
        \item If you answered `yes', give full details (e.g. “I pointed Joe Smith to section 2.3 since he didn’t know how to proceed with Question 2”)\\
        
        \begin{tcolorbox}[fit,height=4cm, width=15cm, blank, borderline={1pt}{-2pt},nobeforeafter]
        
        \begin{soln}
        \end{soln}
        
        \end{tcolorbox}
        
    \end{enumerate}
    
    \item 
    \begin{enumerate}
        \item Did you find or come across code that implements any part of this assignment?
        \begin{soln} Yes / No. \end{soln} 
        
        \item If you answered `yes', give full details (book \& page, URL \& location within the page, etc.).\\
        
        \begin{tcolorbox}[fit,height=4cm, width=15cm, blank, borderline={1pt}{-2pt},nobeforeafter]
        
        \begin{soln}
        \end{soln}
        
        \end{tcolorbox}
        
    \end{enumerate}
\end{enumerate}

\end{document}